% =============================================================================
\section{Matriz de Propiedades de Seguridad del Protocolo y Límites Declarados}
\label{sec:tabla_propiedades_limites}
% =============================================================================

En todo proyecto de investigación en ciberseguridad resulta indispensable establecer con claridad y transparencia cuáles propiedades de seguridad están cubiertas por el diseño propuesto y cuáles quedan fuera del alcance de la implementación actual.

Para organizar esta información de manera comprensible, en la Tabla \ref{tab:propiedades_limites_protocolo} se presenta la relación entre cada propiedad de seguridad electoral estudiada en este capítulo, indicando si se encuentra cubierta por el protocolo, qué tecnología se utiliza para implementarla y qué tipo de prueba respalda su cumplimiento.

\begin{table}[!htbp]
\centering
\renewcommand{\arraystretch}{1.3}
\small
\begin{tabularx}{\textwidth}{|>{\bfseries}p{3.2cm}|c|>{\raggedright\arraybackslash}p{4.2cm}|>{\raggedright\arraybackslash}X|}
\hline
\textbf{Propiedad} & \textbf{¿Cubierta?} & \textbf{Mecanismo Tecnológico} & \textbf{Tipo de Evidencia / Soporte} \\
\hline
Ejecutabilidad del Protocolo & Sí & Secuencia ordenada de pasos de registro, emisión y conteo & Lema formal de alcanzabilidad verificado en Tamarin Prover \\
\hline
Confidencialidad del Voto & Sí (parcial) & Esquema híbrido: ML-KEM-768 + AES-256-GCM con datos asociados & Lema formal de secreto verificado en Tamarin Prover y seguridad IND-CCA2 \\
\hline
Integridad del Sufragio & Sí & Firma digital post-cuántica ML-DSA-65 (FIPS 204) & Lema formal de integridad y seguridad sobre Module-SIS \\
\hline
Autenticidad y Anti-Replay & Sí & Credencial emitida por la AR y validación de firma del votante & Lema formal de autenticidad en Tamarin Prover \\
\hline
Unicidad y Elegibilidad & Sí & Control transaccional en base de datos y consumo de estado lineal & Lema formal de unicidad verificado en Tamarin Prover \\
\hline
Verificabilidad E2E Individual & Parcial & Recibo electrónico firmado con código resumen SHA3-256 & Pruebas funcionales de consulta del recibo mediante hash \\
\hline
Verificabilidad Universal & No & Requeriría un tablón público completo o red de mezcla con pruebas de conocimiento cero & Límite declarado del diseño actual (propuesto como trabajo futuro) \\
\hline
Libertad de Recibo & No & No se implementan credenciales falsas ni cabinas con canales libres de rastro & Límite declarado en los alcances y límites de la investigación (Sección \ref{sec:alcances_limites_tesis}) \\
\hline
Resistencia a la Coacción & No & No se incluye revocación de votos previos ante vigilancia física de un coactor & Límite declarado en los alcances y límites de la investigación (Sección \ref{sec:alcances_limites_tesis}) \\
\hline
\end{tabularx}
\caption{Propiedades de seguridad y límites del protocolo.}
\label{tab:propiedades_limites_protocolo}
\par\vspace{3pt}{\footnotesize\textit{Fuente: Elaboración propia.}}
\end{table}

Tal como se observa en la Tabla \ref{tab:propiedades_limites_protocolo}, el protocolo propuesto garantiza de forma completa la ejecutabilidad, la integridad, la autenticidad, la elegibilidad y la unicidad del voto, así como la confidencialidad del sufragio frente a atacantes en la red y frente a servidores separados. Por otro lado, propiedades adicionales contra la coacción presencial (como la libertad de recibo y la resistencia a la coacción física) se declaran explícitamente como límites fuera del alcance del prototipo, manteniendo coherencia con los objetivos definidos en el Capítulo 1.
